\newpage
\section{\tkzClass{rectangle} 类}

变量 \code{R} 持有一个表，用于存放矩形。它是可选的，你可以自由选择变量名。不过，为清晰和一致起见，推荐使用 \code{R}。如果你使用自定义变量（例如 rectangles），则必须手动初始化它。如果使用了 \code{R}，\code{init\_elements()} 函数会重新初始化该表。


\subsection{矩形的属性}
\label{sub:rectangle_attributes}


各点按正向（逆时针）给出。不做检验：四个点必须构成一个矩形。

\begin{mybox}
创建 | R.ABCD = rectangle:new(z.A, z.B, z.C, z.D)|
\end{mybox}

  \bgroup
  \catcode`_=12
  \small
  \captionof{table}{rectangle 的属性。}\label{rectangle:attributes}
  \begin{tabular}{lll}
  \toprule
  \textbf{属性}       & \textbf{用法} & \\
  \midrule
  \tkzAttr{rectangle}{pa}     & |z.A = R.ABCD.pa| & \\
  \tkzAttr{rectangle}{pb}     & |z.B = R.ABCD.pb| & \\
  \tkzAttr{rectangle}{pc}     & |z.C = R.ABCD.pc| & \\
  \tkzAttr{rectangle}{pd}     & |z.D = R.ABCD.pd| & \\
  \tkzAttr{rectangle}{type}   &  |R.ABCD.type= 'rectangle'|  &\\
  \tkzAttr{rectangle}{center} & |z.I = R.ABCD.center| & 矩形的中心\\
  \tkzAttr{rectangle}{length} &  |R.ABCD.length| & 长 \\
  \tkzAttr{rectangle}{width}  &  |R.ABCD.width| & 宽 \\
  \tkzAttr{rectangle}{diagonal}  &  |R.ABCD.diagonal| & 对角线长度\\
  \tkzAttr{rectangle}{ab}     &  |R.ABCD.ab|   &  过两个顶点的直线   \\
  \tkzAttr{rectangle}{ac}     &  |R.ABCD.ac|   &  同上。 \\
  \tkzAttr{rectangle}{da}     &  |R.ABCD.da|   &  同上。 \\
  \tkzAttr{rectangle}{bc}     &  |R.ABCD.bc|   &  同上。 \\
  \tkzAttr{rectangle}{bd}     &  |R.ABCD.bd|   &  同上。 \\
  \tkzAttr{rectangle}{cd}     &  |R.ABCD.cd|   &  同上。 \\
  \bottomrule
  \end{tabular}
  \egroup



\subsubsection{示例}

\begin{minipage}{.5\textwidth}
\begin{verbatim}
\directlua{
 init_elements()
 z.A = point(0, 0)
 z.B = point(4, 0)
 z.C = point(4, 4)
 z.D = point(0, 4)
 R.new = rectangle(z.A, z.B, z.C, z.D)
 z.I = R.new.center}

\begin{tikzpicture}
\tkzGetNodes
\tkzDrawPolygon(A,B,C,D)
\tkzDrawPoints(A,B,C,D)
\tkzLabelPoints(A,B)
\tkzLabelPoints[above](C,D)
\tkzDrawPoints[red](I)
\end{tikzpicture}
\end{verbatim}
\end{minipage}
\hspace{\fill}\begin{minipage}{.5\textwidth}
\directlua{
 init_elements()
 z.A = point(0, 0)
 z.B = point(4, 0)
 z.C = point(4, 4)
 z.D = point(0, 4)
 R.new = rectangle(z.A, z.B, z.C, z.D)
 z.I = R.new.center}

\begin{tikzpicture}[scale =1.5]
 \tkzGetNodes
 \tkzDrawPolygon(A,B,C,D)
 \tkzDrawSegment[dashed](A,C)
 \tkzDrawPoints(A,B,C,D)
 \tkzLabelPoints(A,B)
 \tkzLabelPoints[above](C,D)
 \tkzDrawPoints[red](I)
 \tkzLabelPoint[right = 10pt](I){$I$\\ |R.new.center|}
 \tkzLabelSegment[sloped,above](C,D){|R.new.length| =
    \pmpn{\tkzUseLua{R.new.length}}}
 \tkzLabelSegment[sloped,above](A,C){|R.new.diagonal| =
   \pmpn{\tkzUseLua{R.new.diagonal}}}
   % \tkzUseLua{R.new.length} and \tkzUseLua{R.new.diagonal} to get the values.
\end{tikzpicture}
\end{minipage}
\newpage

\subsection{矩形的方法}

  \bgroup
  \catcode`_=12
  \small
  \captionof{table}{矩形的方法。}\label{rectangle:methods}
  \begin{tabular}{lll}
  \textbf{方法} & \texttt{参考}  &  \\
  \toprule
  \tkzMeth{rectangle}{new(za ,zb, zc, zd)} &[\ref{meth_rectangle_new}]\\
  \tkzMeth{rectangle}{angle (zi, za, angle)} & [\ref{ssub:angle_method}]\\
  \tkzMeth{rectangle}{gold (za, zb)} & [\ref{ssub:gold_method}]\\
  \tkzMeth{rectangle}{diagonal (za, zc)} &[\ref{ssub:diagonal_method}]\\
  \tkzMeth{rectangle}{side (za, zb, d)} & [\ref{ssub:side_method}]\\
  \tkzMeth{rectangle}{get\_lengths ()} & [\ref{ssub:function_rectangle_get__lengths}]\\
  \tkzMeth{rectangle}{get(<i>)} & [\ref{ssub:method_rectangle_get}]\\
  \bottomrule %
  \end{tabular}
  \egroup


\subsubsection{方法 \tkzMeth{rectangle}{get(<i>)}}
\label{ssub:method_rectangle_get}

该方法返回定义矩形对象的各个点。

\begin{itemize}
  \item \code{R.ABCD:get()} 按顺序返回四个顶点：\code{pa}、\code{pb}、\code{pc}、\code{pd}；
  \item \code{R.ABCD:get(i)}（\code{i} 取 1 到 4）只返回第 \code{i} 个顶点；例如 \code{R.ABCD:get(2)} 等价于 \code{R.ABCD.pb}。
\end{itemize}

\begin{tkzexample}[code only]
\directlua{
  init_elements()
  z.A = point(0, 0)
  z.B = point(2, 0)
  R.ABCD = rectangle:side(z.A, z.B, 1)
  z.E, z.F, z.G, z.H = R.ABCD:get()
  z.K = R.ABCD:get(3)
}
\end{tkzexample}

\subsubsection{方法 \tkzMeth{rectangle}{new(pt,pt,pt,pt)}}
\label{meth_rectangle_new}
该函数用四个点创建一个矩形。不做检验，验证由用户自行负责。

\begin{tkzexample}[latex=.5\textwidth]
  \directlua{
  z.A = point(0, 0)
  z.B = point(4, 0)
  z.C = point(4, 3)
  z.D = point(-2, -3)
  L.AB = line(z.A,z.B)
  L.CD = line(z.C,z.D)
  z.I = intersection(L.AB, L.CD)
  C.I = circle(through(z.I, 3.5))
  z.G,z.E = intersection(L.AB, C.I)
  z.F,z.H = intersection(L.CD, C.I)
  R.I = rectangle(z.E, z.F, z.G, z.H)
  z.X = R.I.ab:projection(z.I)}
  \begin{tikzpicture}
  \tkzGetNodes
  \tkzDrawPolygon(E,F,G,H)
  \tkzDrawPoints(A,B,C,D,I,E,F,G,H,X)
  \tkzLabelPoints(A,B,C,D,I,E,F,G,H,X)
  \tkzDrawPoints(I)
  \end{tikzpicture}
\end{tkzexample}

\subsubsection{方法 \tkzMeth{rectangle}{angle(pt,pt,an)}}
\label{ssub:angle_method}

|R.ang = rectangle:angle(z.I,z.A)|；|z.A| 为顶点；ang 为两个顶点之间的角


\begin{minipage}{.5\textwidth}
\begin{verbatim}
\directlua{
 init_elements()
 z.A = point(0, 0)
 z.B = point(4, 0)
 z.I = point(4, 3)
 P.ABCD = rectangle:angle(z.I, z.A,
          math.pi / 6)
 z.B = P.ABCD.pb
 z.C = P.ABCD.pc
 z.D = P.ABCD.pd}
\begin{tikzpicture}[scale   = .5]
\tkzGetNodes
\tkzDrawPolygon(A,B,C,D)
\tkzDrawPoints(A,B,C)
\tkzLabelPoints(A,B,C,D)
\tkzDrawPoints[new](I)
\end{tikzpicture}
\end{verbatim}
\end{minipage}
\begin{minipage}{.5\textwidth}
\directlua{
 init_elements()
 z.A = point(0, 0)
 z.B = point(4, 0)
 z.I = point(4, 3)
 P.ABCD = rectangle:angle(z.I, z.A, math.pi / 6)
 z.B = P.ABCD.pb
 z.C = P.ABCD.pc
 z.D = P.ABCD.pd}
\begin{tikzpicture}[scale  = .5]
\tkzGetNodes
\tkzDrawPolygon(A,B,C,D)
\tkzDrawPoints(A,B,C)
\tkzLabelPoints(A,B)
\tkzLabelPoints[above](C,D)
\tkzDrawPoints[new](I)
\tkzLabelSegment[sloped,above](A,B){%
    |rectangle: angle (z.C,z.A,math.pi/6)|}
\end{tikzpicture}
\end{minipage}

\subsubsection{方法 \tkzMeth{rectangle}{side(pt,pt,d)}}
\label{ssub:side_method}
\begin{minipage}{.5\textwidth}
\begin{verbatim}
\directlua{
 init_elements()
 z.A = point(0, 0)
 z.B = point(4, 3)
 R.side = rectangle:side(z.A, z.B, 3)
 z.C = R.side.pc
 z.D = R.side.pd
 z.I = R.side.center}
\begin{tikzpicture}
\tkzGetNodes
\tkzDrawPolygon(A,B,C,D)
\tkzDrawPoints(A,B,C,D)
\tkzLabelPoints(A,B)
\tkzLabelPoints[above](C,D)
\tkzDrawPoints[red](I)
\end{tikzpicture}
\end{verbatim}
\end{minipage}
\begin{minipage}{.5\textwidth}
\directlua{
 init_elements()
 z.A = point(0, 0)
 z.B = point(4, 3)
 R.side = rectangle:side(z.A, z.B, 3)
 z.C = R.side.pc
 z.D = R.side.pd
 z.I = R.side.center}
\begin{tikzpicture}
\tkzGetNodes
\tkzDrawPolygon(A,B,C,D)
\tkzDrawPoints(A,B,C,D)
\tkzLabelPoints(A,B)
\tkzLabelPoints[above](C,D)
\tkzDrawPoints[red](I)
\tkzLabelSegment[sloped,above](A,B){%
   |rectangle:side(z.A,z.B,3)|}
\end{tikzpicture}
\end{minipage}

\subsubsection{方法 \tkzMeth{rectangle}{diagonal(pt,pt)}}
\label{ssub:diagonal_method}
矩形的边平行于坐标轴：仅凭一条对角线并不能确定一个矩形。

\begin{minipage}{.5\textwidth}
\begin{verbatim}
\directlua{
init_elements()
z.A = point(0, 0)
z.C = point(4, 3)
R.diag = rectangle:diagonal(z.A, z.C)
z.B = R.diag.pb
z.D = R.diag.pd
z.I = R.diag.center}
\begin{tikzpicture}
\tkzGetNodes
\tkzDrawPolygon(A,B,C,D)
\tkzDrawPoints(A,B,C,D)
\tkzLabelPoints(A,B)
\tkzLabelPoints[above](C,D)
\tkzDrawPoints[red](I)
\tkzLabelSegment[sloped,above](A,B){%
    |rectangle:diagonal(z.A,z.C)|}
\end{tikzpicture}
\end{verbatim}
\end{minipage}
\begin{minipage}{.5\textwidth}
\directlua{
 z.A = point(0, 0)
 z.C = point(4, 3)
 R.diag = rectangle:diagonal(z.A, z.C)
 z.B = R.diag.pb
 z.D = R.diag.pd
 z.I = R.diag.center}

\begin{tikzpicture}
\tkzGetNodes
\tkzDrawPolygon(A,B,C,D)
\tkzDrawPoints(A,B,C,D)
\tkzLabelPoints(A,B)
\tkzLabelPoints[above](C,D)
\tkzDrawPoints[red](I)
\tkzLabelSegment[sloped,above](A,B){%
   |rectangle:diagonal(z.A, z.C)|}
\end{tikzpicture}
\end{minipage}

\subsubsection{方法 \tkzMeth{rectangle}{gold(pt,pt)}}
\label{ssub:gold_method}
\begin{minipage}{.5\textwidth}
\begin{verbatim}
\directlua{
 init_elements()
 z.X = point(0, 0)
 z.Y = point(4, 2)
 R.gold = rectangle:gold(z.X, z.Y)
 z.Z = R.gold.pc
 z.W = R.gold.pd
 z.I = R.gold.center}
\begin{tikzpicture}
\tkzGetNodes
\tkzDrawPolygon(X,Y,Z,W)
\tkzDrawPoints(X,Y,Z,W)
\tkzLabelPoints(X,Y)
\tkzLabelPoints[above](Z,W)
\tkzDrawPoints[red](I)
\tkzLabelSegment[sloped,above](X,Y){%
|rectangle:gold(z.X,z.Y)|}
\end{tikzpicture}
\end{verbatim}
\end{minipage}
\begin{minipage}{.5\textwidth}
\directlua{
init_elements()
z.X = point(0, 0)
z.Y = point(4, 2)
R.gold = rectangle:gold(z.X, z.Y)
z.Z = R.gold.pc
z.W = R.gold.pd
z.I = R.gold.center}

\begin{tikzpicture}
\tkzGetNodes
\tkzDrawPolygon(X,Y,Z,W)
\tkzDrawPoints(X,Y,Z,W)
\tkzLabelPoints(X,Y)
\tkzLabelPoints[above](Z,W)
\tkzDrawPoints[red](I)
\tkzLabelSegment[sloped,above](X,Y){%
  |rectangle:gold(z.X,z.Y)|}
\end{tikzpicture}
\end{minipage}

\subsubsection{方法 \tkzMeth{rectangle}{get\_lengths}()}
\label{ssub:function_rectangle_get__lengths}

\begin{tkzexample}[latex=.45\textwidth]
\directlua{
 init_elements()
 z.I = point(2, 1)
 z.A = point(0, 0)
 R.ABCD = rectangle:angle(z.I, z.A, math.pi / 3)
 z.B = R.ABCD.pb
 z.C = R.ABCD.pc
 z.D = R.ABCD.pd
 tkzx,tkzy = R.ABCD:get_lengths()}
\begin{tikzpicture}
\tkzGetNodes
\tkzDrawPolygon(A,B,C,D)
\tkzDrawCircle(I,A)
\tkzDrawPoints(A,B,C,D)
\tkzLabelPoints(A,B)
\tkzLabelPoints[above](C,D)
\tkzDrawPoints[](I)
\tkzLabelSegment(A,B){%
    $\pmpn{\tkzUseLua{tkzx}}$}
\tkzLabelSegment[right](B,C){%
    $\pmpn{\tkzUseLua{tkzy}}$}
\end{tikzpicture}
\end{tkzexample}
\endinput